%! AP Statistics — Unit 4, Lesson 4.7 — Guided Notes
\documentclass[10pt]{article}
\usepackage{apstats-article}
\usepackage{apstats-boxes}
\newcommand{\TallMath}[1]{$\displaystyle #1\rule[-1.4em]{0pt}{3.2em}$}

\begin{document}

\pageheader{Unit 4, Lesson 4.7}{Guided Notes}
\namedateperiod

\vspace{0.12in}

\begin{objectivebox}
\small By the end of this lesson, I will be able to\dots\\[4pt]
\writeline\\[1pt]
\writeline\\[1pt]
\writeline
\end{objectivebox}

\vspace{0.1in}

\begin{vocabbox}
\termblanklong{Random variable}
\termblanklong{Discrete random variable}
\termblanklong{Probability distribution}
\termblanklong{Probability histogram}
\termblanklong{Cumulative probability distribution}
\end{vocabbox}

\vspace{0.1in}

\begin{hookbox}
\small\textbf{Number of pets owned per household:}
\begin{center}
\renewcommand{\arraystretch}{1.3}
\begin{tabular}{@{}ccccc@{}}
Pets ($x$) & 0 & 1 & 2 & 3 \\
\hline
Probability & 0.30 & 0.40 & 0.20 & 0.10 \\
\end{tabular}
\end{center}
\textbf{What is $P(X \le 1)$?} \quad \blank{4cm} \quad
\textbf{What must the probabilities sum to?} \quad \blank{3cm}
\end{hookbox}

\vspace{0.1in}

\begin{notesbox}{Random Variables and Probability Distributions}
\small

\textbf{Definition:} A \textbf{random variable} is a variable whose value is a
\blank{5cm} of a random process.

\vspace{8pt}
A \textbf{discrete random variable} can take a \blank{4cm} number of values.
Each value has a probability associated with it.

\vspace{10pt}
\textbf{Valid Probability Distribution --- Two Conditions:}
\begin{enumerate}[leftmargin=2em, itemsep=4pt]
  \item Each probability: \blank{6cm}
  \item Sum of all probabilities: \blank{5cm}
\end{enumerate}

\vspace{10pt}
\textbf{Example:} $X$ = number of siblings for a randomly chosen student.

\vspace{4pt}
\renewcommand{\arraystretch}{1.7}
\begin{tabularx}{\linewidth}{@{} l X X X X X @{}}
$x$ & 0 & 1 & 2 & 3 & 4 \\
\hline
$P(X=x)$ & \blank{1.3cm} & \blank{1.3cm} & \blank{1.3cm} & \blank{1.3cm} & \blank{1.3cm} \\
\end{tabularx}

\vspace{8pt}
Check conditions: \quad (1) \blank{5cm} \quad (2) \blank{5cm}

\end{notesbox}

\vspace{0.1in}

\begin{notesbox}{Cumulative Probability Distribution}
\small

$P(X \le x)$ = the probability that $X$ is \blank{5cm}.

\vspace{8pt}
\textbf{Example:} Use the siblings distribution above.

\renewcommand{\arraystretch}{1.7}
\begin{tabularx}{\linewidth}{@{} l X X X X X @{}}
$x$ & 0 & 1 & 2 & 3 & 4 \\
\hline
$P(X=x)$ & 0.15 & 0.40 & 0.30 & 0.10 & 0.05 \\
$P(X \le x)$ & \blank{1.3cm} & \blank{1.3cm} & \blank{1.3cm} & \blank{1.3cm} & \blank{1.3cm} \\
\end{tabularx}

\vspace{8pt}
$P(X \ge 3) = $ \blank{4cm} (hint: use complement or add $P(3)$ and $P(4)$)

\vspace{8pt}
Interpret $P(X \le 2) = 0.85$ in context:\\
\writeline
\end{notesbox}

\vspace{0.1in}

\begin{notesbox}{Interpreting a Probability Distribution}
\small

When you describe a probability distribution, discuss \textbf{three things}:

\vspace{6pt}
\renewcommand{\arraystretch}{1.8}
\begin{tabularx}{\linewidth}{@{} p{2.2cm} X @{}}
\textbf{Shape} & \blank{10cm} \\
\textbf{Center} & Where are the probabilities \blank{3cm}? \\
\textbf{Spread} & How \blank{5cm} are the values? \\
\end{tabularx}

\vspace{8pt}
\textbf{Interpretation of the siblings distribution in context:}\\
\writeline\\[1pt]\writeline
\end{notesbox}

\vspace{0.1in}

\begin{practicebox}
\small
$X$ = number of cars owned per household. Use the distribution:
\begin{center}
\renewcommand{\arraystretch}{1.3}
\begin{tabular}{@{}ccccc@{}}
$x$ & 0 & 1 & 2 & 3 \\
\hline
$P(X=x)$ & 0.10 & 0.40 & 0.35 & 0.15 \\
\end{tabular}
\end{center}
\begin{enumerate}[leftmargin=1.5em, itemsep=8pt]
  \item Verify this is a valid probability distribution: \blank{6cm}
  \item $P(X = 2) = $ \blank{2cm}
  \item $P(X \le 1) = $ \blank{2cm}
  \item $P(X \ge 2) = $ \blank{2cm}
  \item Interpret $P(X \ge 2)$ in context:\\
        \writeline
\end{enumerate}
\end{practicebox}

\end{document}
